\documentclass[a4paper,12pt]{article}

% Import the deliverable package from common directory
\usepackage{../common/deliverable}

% Tell LaTeX where to find graphics files
\graphicspath{{../common/logos/}{./figures/}{../}}

\usepackage{xspace}
\usepackage{lipsum}

\usepackage{amsmath}
\usepackage{bm}
\usepackage{subcaption}
\usepackage{algorithm}
\usepackage{algorithmic}
\usepackage{multirow}


\usepackage{tikz}
\usepackage{pgfplots}
\usepackage{pgfplotstable}
\usepackage{listings}
\lstset{
    language=C++,
    basicstyle=\ttfamily\scriptsize,
    keywordstyle=\color{blue},
    commentstyle=\color{gray},
    stringstyle=\color{red},
    breaklines=true,
    frame=single,
    numbers=left,
    numberstyle=\tiny,
    showstringspaces=false
}

\pgfplotsset{compat=1.9}

% Set the deliverable number (without the D prefix, it's added automatically)
\setdeliverableNumber{1.9}

% Begin document
\begin{document}

% Create the title page with the title as argument
\maketitlepage{Report on deal.II improvements -- IV semester}

\newpage

% Main Table using the new environment and command
\begin{deliverableTable}
    \tableEntry{Deliverable title}{Report on deal.II improvements -- IV semester}
    \tableEntry{Deliverable number}{D1.9}
    \tableEntry{Deliverable version}{v1.1}
    \tableEntry{Date of delivery}{July 31, 2026}
    \tableEntry{Actual date of delivery}{July 31, 2026}
    \tableEntry{Nature of deliverable}{Report}
    \tableEntry{Dissemination level}{Public}
    \tableEntry{Work Package}{WP1}
    \tableEntry{Partner responsible}{INPT, RUB}
\end{deliverableTable}

% Abstract and Keywords Section
\begin{deliverableTable}
    \tableEntry{Abstract}{This deliverable reports the improvements made to the deal.II
    	finite element library during the fourth semester of the dealii-X project. It covers
    	the five activities of Work Package 1: the extension of the exascale capabilities of
    	deal.II on GPU accelerators, including batched matrix-free kernels, Raviart--Thomas
    	elements, discontinuous Galerkin face integration and a portable multigrid
    	infrastructure (RUB); an implicit hybridizable discontinuous Galerkin solver for
    	one-dimensional blood flow in vascular networks (UNIPI); the polygonal and polyhedral
    	discretization module and its submission to the deal.II core library, together with an
    	unfitted space-time formulation on moving domains (SISSA); the integration of the
    	PSCToolkit algebraic multigrid preconditioners and their assessment on MareNostrum 5
    	(UNITOV); and the integration of the MUMPS sparse direct solver, with block low-rank
    	compression and mixed-precision iterative refinement applied to organ-scale problems
    	(INPT). For each activity the report describes the implementation, the software
    	artefacts released, and the performance obtained on European pre-exascale systems.}
    \tableEntry{Keywords}{deal.II; exascale computing; matrix-free GPU kernels; multigrid
    	and algebraic preconditioners; sparse direct solvers}
\end{deliverableTable}

\newpage

\begin{documentControl}
    \addVersion{0.1}{July 18, 2026}{Biswajeet Rath}{Initial draft}
    \addVersion{0.2}{July 18, 2026}{Ivan Prusak}{Restructuring and adjustments}
    \addVersion{0.3}{July 30, 2026}{Marco Feder}{Update with PSCToolkit}
    \addVersion{0.4}{July 31, 2026}{Luca Heltai}{Refinement on UNIPI part}
    \addVersion{0.5}{July 31, 2026}{Ivan Prusak}{Final review and adjustments}
    \addVersion{1.0}{July 31, 2026}{Luca Heltai}{Final version for submission}
    \addVersion{1.1}{September 3, 2026}{Luca Heltai}{Revision following the intermediate review}
\end{documentControl}

\subsection*{{Approval Details}}
Approved by: Martin Kronbichler \\
Approval Date: 31.07.2026

\subsection*{{Distribution List}}
\begin{itemize}
    \item [] - Project Coordinators (PCs)
    \item [] - Work Package Leaders (WPLs)
    \item [] - Steering Committee (SC)
    \item [] - European Commission (EC)
\end{itemize}

\vspace*{2cm}

\disclaimer

\newpage

\tableofcontents % Automatically generated and hyperlinked Table of Contents

\newpage

\section{{Introduction}}

The dealii-X project is extending the \texttt{deal.II} finite element library and its surrounding solver ecosystem with numerical methods, software abstractions, and high-performance implementations required by exascale simulations of digital twins of the human body. Reaching this objective requires more than increasing the size of existing computations. It also requires reducing data movement, exposing sufficient parallelism on heterogeneous architectures, improving the scalability and robustness of linear and nonlinear solvers, and introducing modelling strategies that preserve the relevant physical information at substantially lower computational cost.

\subsection{{Purpose of the Document}}
This deliverable reports the improvements developed during the fourth project seme\-ster. The contributions address complementary layers of the simulation stack. UNIPI develops geometrically reduced flow models on metric graphs and an implicit HDG implementation for network haemodynamics; INPT advances block low-rank and mixed-precision algorithms in MUMPS; SISSA develops polygonal and unfitted discretization technology; RUB improves performance-portable matrix-free kernels and multigrid infrastructure; and UNITOV completes the integration of PSBLAS into the \texttt{deal.II} library, advances the AMG4PSBLAS interface, and evaluates the resulting algebraic multigrid solvers on MareNostrum~5. The UNITOV contribution is particularly important because it turns PSCToolkit from an external capability into a native \texttt{deal.II} linear-algebra backend with familiar vector, matrix, solver, and preconditioner interfaces, thereby making scalable sparse algebra directly available to application developers.

A common theme of the reported work is the use of mathematical and computational structure to reduce time-to-solution. Depending on the problem, this structure may be algebraic, geometric, induced by the discretization, or exposed by the hardware. Algebraic low-rank approximations compress matrix blocks whose action can be represented accurately in a lower-dimensional subspace. Geometrical model reduction replaces selected three-dimensional network-like subdomains by lower-dimensional representations when the dominant dynamics are constrained by their slender geometry. Polygonal and unfitted methods increase geometric flexibility, matrix-free methods reduce memory traffic, and scalable sparse algebra supplies robust solvers for problems for which explicit matrices remain necessary. Although these approaches operate at different levels, they pursue the same objective: representing and solving the relevant high-dimensional problem with fewer degrees of freedom, less data movement, or fewer arithmetic and communication operations, while retaining controlled accuracy.




\subsection{{Structure of the Document}}
\begin{itemize}
    \item Section~\ref{sec:rub}: WP1.1: Extending and improving the exascale capabilities of deal.II (RUB)
    \item Section~\ref{sec:unipi}: WP1.2: Improvement of pre-exascale modules of the deal.II library (UNIPI)
    \item Section~\ref{sec:sissa}: WP1.3: Experimental polygonal discretization module for deal.II (SISSA)
    \item Section~\ref{sec:unitov}: WP1.4: Integration of PSCToolkit within deal.II (UNITOV)
    \item Section~\ref{sec:inpt}: WP1.5: Integration of MUMPS within deal.II (INPT)
\end{itemize}


\newpage


\section{WP1.1: Extending and improving the exascale capabilities of deal.II (RUB)}
\label{sec:rub}

\subsection{{Performance Evaluations of High Order Finite Element Kernels on Modern Accelerators}}

We detail the implementation techniques and performance measurements of high-order matrix-free finite element methods on graphics processing units (GPUs).
The matrix-free framework and its initial implementations are described in Deliverable D1.3
{\footnotesize \url{https://github.com/dealii-X/deliverables/blob/main/D1.3/D1.3_dealii_report.pdf}.\\}
To evaluate these problems, we adopt the bake-off benchmarks \citep{bakeoff} as our baseline problems. In this report, we also implement a new technique to improve operator performance for low-order elements.
Previous high-order matrix-free GPU implementations were based on assigning a single element per thread block \citep{tens_prod}. This approach causes low-order elements to underutilize
the streaming multiprocessors (SMs) because NVIDIA only allows each SM to schedule 32 thread blocks \citep{nvidia_hopper_tuning}. Consequently, low-order elements reach this limit,
leaving registers and shared memory underutilized within each SM. On AMD GPUs, the single element per thread block approach leaves the 64-thread wavefront underutilized for low-order elements, consequently causing poor performance.
Based on a batch processing approach, our technique assigns multiple elements to each thread block to overcome these limitations. For example, a polynomial order of 1 in the Laplace operator evaluation (Bake-Off Kernel 3 in the CEED benchmarks) can fit 11 elements per thread block.
A performance comparison between the single-element-per-thread-block and the batched processing approaches is given in Figure~\ref{fig:BK3_1_10}. The new technique achieves a 2$\times$ speedup for polynomial order 1 on the NVIDIA GH200 Superchip at the JUPITER supercomputer located at the Jülich Research Centre.
New batched implementations are available at \url{https://github.com/dealii-X/benchmarks/tree/main/CEED_BK}.\footnote{retrieved on 28.07.2026}
\\

\begin{figure}[htbp]
	\centering
	\includegraphics[width=0.9\textwidth]{figures/GH200_fp64_BK3_1KB_10KB}
	\caption{Performance comparison of the Laplace operator (BK3) in FP64 on the NVIDIA GH200. Single-element-per-thread-block approach (left); proposed multiple-elements-per-thread-block approach (right).}
	\label{fig:BK3_1_10}
\end{figure}

We refer to Deliverable D3.4 for the detailed performance analysis with NVIDIA Nsight Compute as well as the roofline model.


\subsection{Laplace Solver Performance (CEED BP3)}
\begin{figure}[htbp]
	\centering
	\includegraphics[width=0.5\textwidth]{figures/GH200_BK_BP}
	\caption{
		Performance results for the Laplace Operator (CEED BK3) \& Problem (CEED BP3) are evaluated across various polynomial orders and problem sizes on the GH200 GPU. Both Bake-Off Kernel performance (GDOF/s)
		and Bake-Off Problem performance ($\text{GDOF} \times \text{CG}_{\text{iter}} / (\text{second} \times \text{process})$) are measured using the global degrees of freedom (the T-vector in CEED).
	}
	\label{fig:BK_BP}
\end{figure}

We have extended our operator evaluation to a complete linear solver by embedding the Laplace operator within a Conjugate Gradient (CG) method, corresponding to the CEED Bake-Off Problem 3 (BP3) benchmark.
The solver, implemented within the \texttt{deal.II} finite element library framework, is evaluated on both single and multi-GPU configurations of the NVIDIA GH200 Superchip, utilizing MPI with peer-to-peer (P2P) communication enabled. Figure~\ref{fig:BK_BP} illustrates the initial performance results,
highlighting that the performance difference between operator evaluations (BK) and full linear solvers (BP)
is primarily driven by gather-scatter overheads and inter-GPU communication. Additionally, the results confirm the computational efficiency of high-order elements which directly leverage the tensor-product formulation of the shape functions to increase arithmetic intensity.
The source code can be found at \url{https://github.com/dealii-X/benchmarks/tree/main/CEED_bp}.\footnote{retrieved on 28.07.2026}



\subsection{Raviart--Thomas Elements and Kernel Performance}

\begin{figure}[htbp]
	\centering
	\includegraphics[width=0.9\textwidth]{figures/GH200_Mass_32_64}
	\caption{Performance comparison of the vector Mass operator using FP32 (left) and FP64 (right) precisions on the NVIDIA GH200 Superchip.}
	\label{fig:GH200_vector_Mass}
\end{figure}

\begin{figure}[htbp]
	\centering
	\includegraphics[width=0.9\textwidth]{figures/MI250x_Mass_32_64}
	\caption{Performance comparison of the vector Mass operator using FP32 (left) and FP64 (right) precisions on the single AMD MI250X GPU tile.}
	\label{fig:MI250x_vector_Mass}
\end{figure}
$H^1$-conforming elements (as used in the CEED Bake-Off kernels) enforce continuity of all vector components across elements, whereas the Navier-Stokes equations
in mixed formulations often require continuity of the normal component of the flux across element boundaries which is a property of $H(\mathrm{div})$-conforming spaces.
One element type designed for this purpose are the Raviart--Thomas finite elements. We developed the vector Mass, mixed divergence and mixed gradient operators, which are the core operations
in Raviart-Thomas kernel design, and implemented them in our benchmarks. These benchmarks are available at \url{https://github.com/dealii-X/benchmarks/tree/main/rt\_kernels}.\footnote{retrieved on 28.07.2026}




Figure \ref{fig:GH200_vector_Mass} shows the performance of the vector Mass operator on the NVIDIA GH200 in both FP32 and FP64 precisions. Batch processing enabled all polynomial orders to achieve similar throughput (GDOF/s). Similar trends are observed for the AMD MI250X
GPU om the LUMI supercomputer, which performs well given its $1.3\text{ TB/s}$ STREAM bandwidth. Only polynomial order 8 performs significantly worse than the other orders, as shown in Figure \ref{fig:MI250x_vector_Mass}.
This drop in performance is likely due to register file capacity constraints on this slightly older AMD architecture.


We also introduce a new technique to target NVIDIA tensor cores \citep{nvidia_mma}
which are designed to perform fast matrix-matrix multiplication in different precisions. Initial implementations for FP64 and non-IEEE-compliant TensorFloat-32 (TF32) are available
in the same \texttt{rt\_kernels} benchmarks repository.


\newcommand{\average}[1]{\left\{\!\left\{#1\right\}\!\right\}}
\newcommand{\jump}[1]{\left[\!\left[{#1}\right]\!\right]}


\subsection{Discontinuous Galerkin kernels and multigrid solvers}

Another contribution by the RUB group is the developement of Discontinuous Galer\-kin (DG) kernels including the face integration. In particular, the main interest lies in the matrix-free evaluation of, for instance,  the symmetric interior penalty (Nitsche) discretization of the Laplace operator which reads as
\[
\sum\limits_{K \in \text{cells}} \left(\nabla \varphi_i, \nabla u_h \right)_{K} + \sum\limits_{F \in \text{faces}} \big( -\left< \jump{\varphi_i}, \average{\nabla u_h} \right>_{F} - \left< \average{\nabla \varphi_i}, \jump{u_h}\right>_F  + \left< \jump{\varphi_i}, \sigma \jump{u_h}\right>_F \big),
\]
where $\average{v} = \frac{v^{-}+v^{+}}{2}$ and $\jump{v} = n^{-}v^{-}+n^{+}v^{+}$ are the average and the directed jump of a quantity $v$ over a face between two adjacent cells $K^{-}$ and $K^{+}$, respectively, and $\sigma$ is an appropriate Nitsche penalty parameter to ensure the coercivity of the discretization.



As a first approach, the face-centric loop (FCL) with global temporary storage has been implemented~\citep{KronbichlerAllalen2018, KronbichlerKormann2019}. The FCL algorithm is described as follows:

\noindent\rule{\textwidth}{1pt}\\[-0.5em]
\textbf{FCL algorithm:} Face-centric loop (FCL) with global temporary storage.\\[-0.5em]
\rule{\textwidth}{0.6pt}

\begin{enumerate}[label=\arabic*., leftmargin=2em, itemsep=0.8em]
	\item Parallel for loop over all cells $K$:
	\begin{itemize}[label=--, leftmargin=2em, itemsep=0.2em]
		\item Perform cell integration for all test functions $\varphi_i$ on values of input vector with sum factorization.
		\item Write result into position $i$ of result vector with non-temporal stores.
		\item Interpolate cell values from input vector to all $2d$ faces for both values and the reference-cell normal derivative; write these $2(k+1)^{d-1}$ values per face into global auxiliary face storage with non-temporal stores.
	\end{itemize}
	
	\item Parallel for loop over all inner faces $F$:
	\begin{itemize}[label=--, leftmargin=2em, itemsep=0.2em]
		\item Load values and normal derivatives from global auxiliary storage on elements $K^-$ and $K^+$.
		\item Perform face integration using values and normal derivative with sum factorization, producing results for both sides simultaneously.
		\item Write result tested by value and normal derivative of test function $\varphi_i$ into global auxiliary storage.
	\end{itemize}
	
	\item Parallel for loop over boundary faces $F$: Similar to inner faces.
	
	\item Parallel for loop over all cells $K$:
	\begin{itemize}[label=--, leftmargin=2em, itemsep=0.2em]
		\item Read values and normal derivatives on each face from global auxiliary storage and finalize face integration step by expanding test functions $\varphi_i$ into the cell and adding the contribution into the result vector.
	\end{itemize}
\end{enumerate}

\vspace{-0.5em}
\noindent\rule{\textwidth}{1pt}

\input{figures/dg_matvec.tex}

The throughput performance of the FCL algorithm during matrix-vector product evaluation is presented in Figures~\ref{fig:dg_mat_vec_dp} and~\ref{fig:dg_mat_vec_sp}, alongside the continuous Galerkin operator evaluation. The Degrees of Freedom (DoF) per GPU are measured for each dimension of the corresponding discontinuous and continuous finite element spaces. Continuous matrix–vector evaluation listings are included to complement the DG results rather than to serve as a direct comparison.

Nevertheless, two key differences are worth highlighting: the DG elements reach a clear performance plateau, and increasing the GPU count consistently achieves this same throughput limit. In contrast, the continuous Galerkin elements display a pronounced latency effect. This highlights strong potential for improved scaling—even though the current implementation transmits all ghost cell values (instead of faces) and results in redundant computations across ghost cells shared by MPI processes. Communication reduction and alternative paradigms, such as a cell-centric loop for face integration, are currently under investigation, with a full scalability study planned for the coming months.

\input{figures/dg_mg_jupiter.tex}

To evaluate the performance of the entire solver chain on GPUs, we consider multigrid in the cph-MG variant, which is the most performant one on CPUs: the discontinuous space is first embedded into the corresponding continuous space on the finest level, and the resulting problem is then treated by the ph-multigrid solver developed previously for continuous elements~\citep{Fehn2020}. The performance of the mixed-precision cph-MG solver is shown in Figure~\ref{fig:cph_mg}.


\paragraph{Analysis of the multigrid scalability limitation.}
A complementary weak-scaling campaign on JUPITER, reported in Deliverable
D1.7 and summarized in Deliverable D5.3, separates the scaling of the
fine-grid matrix-free operator from that of the complete multigrid cycle.
With $57$--$65$ million unknowns per GPU, the matrix-vector product retained
approximately $65\%$ parallel efficiency up to 128 nodes (512 GPUs), whereas
the runtime of the complete multigrid method increased by a factor of $2.5$
over the corresponding $128\times$ increase in problem and machine size.
Profiling traced this gap to communication latency on the coarser levels of
the multigrid hierarchy. On these levels, the local computational work
decreases substantially while ghost exchanges and synchronization remain
necessary. The approximately fixed overhead of at least
$7\times10^{-4}\,\mathrm{s}$ per step for these small-message exchanges
therefore dominates an increasing fraction of the multigrid cycle at large
node counts.

The observed degradation is consequently not an unidentified general
scalability limitation of the matrix-free operator evaluation, but a specific
communication bottleneck in the coarse-grid phase of the multigrid algorithm.
Current optimization work focuses on reducing the number and cost of these
communication steps, improving communication--computation overlap, and
investigating alternative coarse-level execution and repartitioning
strategies.

Following the recommendation of the intermediate project review, the
consortium secured support from the POP3 (Performance Optimisation and
Productivity) Centre of Excellence for an independent assessment of this
representative JUPITER workload. The assessment complements the profiling and
co-design activities in WP3 and is intended to provide an independent
characterization of the communication bottleneck and concrete optimization
recommendations. At the time of this revision, the analysis report is
expected in September 2026. In parallel, a researcher with extensive
experience in MPI and communication-avoiding methods has joined the project
for its remaining duration.

The implementation of the algorithms is available at \\ \url{https://github.com/dealii-X/portable-multigrid}.\footnote{retrieved on 28.07.2026}

\subsection{Portable multigrid infrastructure in deal.II}

Developments of the portable multigrid infrastructure is gradually migrating to \texttt{deal.II}. Matrix-free geometric transfer operators were introduced in PR~\href{https://github.com/dealii/dealii/pull/19347}{\#19347}\footnote{retrieved on 28.07.2026} (opened in February 2026 and merged in June 2026). This functionality was subsequently extended to support systems of finite elements in PR~\href{https://github.com/dealii/dealii/pull/20023}{\#20023}\footnote{retrieved on 28.07.2026} (opened and merged in July 2026). Both pull requests are included in the upcoming \texttt{deal.II} release 9.8, scheduled for early August 2026. Furthermore, the latter PR has been successfully utilized in the new tutorial Step-104,\footnote{\url{https://github.com/dealii/dealii/pull/20026}, retrieved on 28.07.2026} demonstrating a substantial speedup of the transfer operators and boosting the overall solver speed by a factor of 3.

\newpage


\section{WP1.2: Improvement of pre-exascale modules of the deal.II library (UNIPI)}
\label{sec:unipi}
%\lipsum[4-5]
Many physiological flow domains are slender and organized as branching networks. Blood vessels and conducting airways have a small cross-section compared with their axial length, while the complete vascular or bronchial tree may contain a very large number of interconnected branches. Resolving every branch with a fully three-dimensional mesh is therefore often unnecessary and, at organ scale, computationally prohibitive. Away from localized regions with strong secondary flow or complex geometric detail, the dominant dynamics can be represented by quantities averaged over each cross-section and evolved along the branch centreline.

This motivates a geometrical reduced-order modelling strategy in which a three-dimensional tubular network is replaced by a one-dimensional metric graph. Graph edges represent individual vessels or airways, while graph vertices represent bifurcations, junctions, inlets, and outlets. The reduced equations retain axial wave propagation, conservation across junctions, compliant-wall effects, and reduced terminal impedances, while eliminating the transverse degrees of freedom associated with a complete three-dimensional discretization. This class of models is directly relevant to systemic and microvascular blood flow and to airway-flow models for the lungs, where a 3D--1D--0D hierarchy can retain detailed three-dimensional descriptions in selected regions and use graph and lumped-parameter models for the many unresolved downstream branches.

Geometrical reduction is complementary to projection-based reduced-order modelling. In a local axial--transverse parametrization, a full-dimensional field may be viewed formally as
\[
u(s,\boldsymbol{\xi},t) \simeq \sum_{k=1}^{r} u_k(s,t)\,\phi_k(\boldsymbol{\xi};s),
\]
where only a small number of transverse modes is retained. The standard one-dimensional model corresponds to an extremely low-rank representation in the cross-sectional variables, usually dominated by averaged area, pressure, and axial velocity or flow rate. Geometrical model reduction can therefore be interpreted as a structured low-rank approximation of the continuous problem: unlike a purely algebraic factorization, the low-rank structure is selected before discretization from geometric slenderness and physical conservation laws. This interpretation connects the UNIPI activity to the broader approximate-computing strategy of dealii-X, in which both geometric and algebraic low-rank approximations replace expensive high-dimensional representations by smaller models while preserving the relevant input--output behaviour.

Within this framework, UNIPI has developed an implicit hybridizable discontinuous Galerkin (HDG) solver for one-dimensional flow on complex metric graphs using \texttt{deal.II}. Blood flow is the principal validation problem, but the core infrastructure---edge-local discretization, numerical traces, conservative junction conditions, implicit time integration, and reduced terminal models---is reusable for other conservation laws on physiological networks, including reduced airway-flow models. The implementation combines Harten--Lax--van Leer (HLL) numerical fluxes with the \textsc{SUNDIALS} IDA solver and uses a fully coupled monolithic Newton method, avoiding nested nonlinear solves at graph junctions and terminal models.

\subsection{Open-source software delivery}

A major result of the UNIPI activity is not only the formulation and validation of the numerical method, but its delivery as openly available and reusable research software. The implementation has been released in the public GitHub repository \url{https://github.com/luca-heltai/metric-flow-x} under the permissive MIT license. The repository contains the C++ implementation in \texttt{source/} and \texttt{include/}, executable drivers in \texttt{apps/}, parameter files, unit and regression tests, mathematical documentation, and benchmark material. Its documented build uses CMake and \texttt{deal.II}. The public release is therefore an explicit software deliverable of the UNIPI contribution, enabling independent inspection, reproduction of the reported results, reuse in other mixed-dimensional applications, and continued community development.

As of 31 July 2026, the repository is public and actively maintained. Its development history returns more than 100 commits, with recent work covering the HDG formulation, the transition to \textsc{SUNDIALS} IDA, analytical-Jacobian verification, MPI compatibility, benchmark networks with 37 and 56 arteries, continuous integration, and user/developer documentation.

\subsection{Governing System}
%\lipsum[6]
We consider the one-dimensional blood flow equations parameterized along a centerline tangent vector $\bm b(\bm x)$. The compact vector form of the conservation system is given by:
\begin{equation}
    \partial_t \mathbf{w} + \nabla \cdot \left(\bm b(\bm{x})  \mathbf{F}(\mathbf{w})\right) = \mathbf{S}(\mathbf{w}), \qquad
    \mathbf{w} = \begin{pmatrix} A \\ U \end{pmatrix}, \qquad
    \mathbf{S}(\mathbf{w}) = \begin{pmatrix} 0 \\ - \frac{f U}{A} \end{pmatrix},
\end{equation}
where $A(x,t)$ is the cross-sectional area, $U(x,t)$ is the mean velocity, $\rho$ is the blood density, and
$f = 2(\zeta+2)\,\mu\,\pi$ is the viscous friction coefficient
with dynamic viscosity $\mu$ and velocity-profile parameter
$\zeta$.
The non-linear physical flux vector $\mathbf{F}(\mathbf{w})$ is:
\begin{equation}
    \mathbf{F}(\mathbf{w}) = \begin{pmatrix} F_A(\mathbf{w}) \\ F_U(\mathbf{w}) \end{pmatrix} = \begin{pmatrix} A U \\ \frac{U^2}{2} + \frac{1}{\rho} P(A) \end{pmatrix}.
\end{equation}
The arterial wall mechanics are governed by the tube law
\begin{align*}
    P(A,x) - P_0 = P_d + \frac{\beta}{A_d(x)}
    \left[ (A^m - A_d(x)^m) \right], \text{ with }  \beta = 4/3 \sqrt{\pi} E h_w(x),
\end{align*}
where $P_0 $ is the external pressure, $P_d$ is the diastolic pressure,
$E$ is the wall Young's modulus,
$h_w(x)$ is the wall thickness, and
$A_d(x)$ is the \emph{local diastolic area}, tube law constant $m =0.5$.

\subsection{HDG Implementation}
The main components implemented in the HDG solver are summarized below:
\begin{itemize}
    \item Implemented the Hybridizable Discontinuous Galerkin (HDG) discretization for the one-dimensional blood flow equations.
    \item Implemented the Harten--Lax--van Leer (HLL) numerical flux for communication between neighboring elements.
    \item Developed a unified trace formulation for
          \begin{itemize}
              \item interior interfaces,
              \item inflow and outlet boundary conditions, and
              \item vascular junctions (including bifurcations).
          \end{itemize}

    \item Incorporated physiological outlet models, including single-resistance and three-element Windkessel (RCR) boundary conditions.
    \item Coupled the semi-discrete HDG system with the \textsc{SUNDIALS} IDA solver for fully implicit time integration.
    \item Assembled the complete residual and analytical Jacobian required by Newton's method for efficient nonlinear solution.
\end{itemize}


\subsection{Results}
%\lipsum[7]
The developed HDG solver has been validated using the well-established 37-artery human systemic network benchmark proposed in~\cite{matthys2007pulse}. The network consists of 37 arterial segments connected through 21 bifurcations and is driven by a physiological inflow waveform. At each terminal vessel, a single-resistance outlet condition,
\[
    P-P_{\mathrm{out}} = RQ,
\]
is imposed, while the bifurcations are coupled through the proposed HDG trace formulation by enforcing mass conservation and continuity of total pressure.

\begin{figure}[ht]
    \centering
    \includegraphics[width=0.45\linewidth]{figures/pressure_at_t14.155_37_arteries.png}
    \caption{Pressure distribution in the 37-artery benchmark at $t\approx14.155\,\mathrm{s}$.}
    \label{fig:37network}
\end{figure}

The numerical solution is verified against the finite-volume benchmark of \cite{boileau2015benchmark}. For comparison, the pressure, flow rate, and cross-sectional area are monitored at the midpoint of the thoracic aorta II, which is one of the major arteries close to the heart. Figure~\ref{fig:thoracic_results} shows an excellent agreement of the computed solution with the benchmark reference throughout one cardiac cycle.

\begin{figure}[ht]
    \centering
    \begin{subfigure}[b]{0.32\textwidth}
        \centering
        \includegraphics[width=\linewidth]{figures/thoracic_aorta_II_pressure.pdf}
        \caption{Pressure}
    \end{subfigure}
    \hfill
    \begin{subfigure}[b]{0.32\textwidth}
        \centering
        \includegraphics[width=\linewidth]{figures/thoracic_aorta_II_flow.pdf}
        \caption{Flow rate}
    \end{subfigure}
    \hfill
    \begin{subfigure}[b]{0.32\textwidth}
        \centering
        \includegraphics[width=\linewidth]{figures/thoracic_aorta_II_area.pdf}
        \caption{Area}
    \end{subfigure}
    \caption{Comparison of pressure, flow rate, and cross-sectional area at the midpoint of the thoracic aorta II over one cardiac cycle.}
    \label{fig:thoracic_results}
\end{figure}

To further demonstrate that the numerical solution is independent of the spatial discretization, a mesh independence study was carried out using successive mesh refinements. As shown in Figure~\ref{fig:mesh_independence}, the numerical solution converges rapidly, with negligible differences observed for the finer meshes. This confirms that the reported results are not artifacts of the computational mesh.

\begin{figure}[ht]
    \centering
    \includegraphics[width=\linewidth]{figures/results_37_thoracic_aorta_II_mesh_independence.png}
    \caption{Mesh independence study for the thoracic aorta II.}
    \label{fig:mesh_independence}
\end{figure}


\newpage


\section{WP1.3: Experimental polygonal discretization module for deal.II (SISSA)}
\label{sec:sissa}

\subsection{Polygonal Discretization Module}

The polygonal discretization activity defined in Work Package 1.3 aims to integrate general-mesh discretization capabilities into \texttt{deal.II}, with particular emphasis on discontinuous Galerkin methods on agglomerated polygonal and polyhedral meshes and on the automatic construction of coarse polygonal grid hierarchies. Deliverable D1.5 introduced the external \texttt{polyDEAL} library and described a Poisson tutorial on agglomerated polytopic meshes.
Deliverable D1.7 subsequently reported the testing and validation of the library, the opening of a separate, self-contained Code Gallery contribution through pull request \href{https://github.com/dealii/code-gallery/pull/233}{\#233},\footnote{retrieved on 28.07.2026} and the ongoing refactoring of reusable components in preparation for their integration into \texttt{deal.II}.
During the present reporting period, the Code Gallery contribution was reviewed and merged, and the complete \texttt{agglomeration\_poisson} program, its documentation, and its numerical results became publicly available on the official \href{https://dealii.org/developer/doxygen/deal.II/code_gallery_agglomeration_poisson.html}{\texttt{deal.II} Code Gallery page}.\footnote{retrieved on 28.07.2026} A follow-up documentation update improving figure rendering and page layout was subsequently submitted through pull request \href{https://github.com/dealii/code-gallery/pull/244}{\#244}.\footnote{retrieved on 28.07.2026}

The public
\href{https://github.com/fdrmrc/Polydeal}
{\texttt{polyDEAL} repository}\footnote{retrieved on 31.07.2026}
also contains an example program demonstrating multigrid on nested
agglomerated meshes. The program constructs an R-tree-based hierarchy of
agglomerated meshes and the associated interlevel transfer operators, and
applies the resulting multigrid method to the fine-grid linear system.

In parallel, the integration activity announced in D1.7 has progressed to a concrete submission to the \texttt{deal.II} core library. The central reusable components originating from \texttt{polyDEAL} were separated from application-specific code, adapted to \texttt{deal.II} interfaces and coding conventions, and extended into a framework for DG discretizations on agglomerated meshes.
The resulting \texttt{deal.II} pull request \href{https://github.com/dealii/dealii/pull/19891}{\#19891}\footnote{retrieved on 28.07.2026} includes agglomeration management, \texttt{deal.II}-style accessor and iterator interfaces, agglomerated degree-of-freedom and sparsity-pattern construction, agglomeration-aware mapping and quadrature, and R-tree-based hierarchical agglomeration. At the reporting date, the pull request is open and under review.
A detailed description of this integration activity is provided in Deliverable D1.8, Polygonal discretization in deal.II -- II.

\subsection{Unfitted space--time DG method on moving domains}

Building on the agglomeration-based polygonal discretization framework developed in Work Package 1.3, related theoretical work has extended the same general-mesh DG philosophy to parabolic problems posed on moving domains. The evolving physical domain is embedded in a fixed background space--time mesh (for example, a rectangular mesh) which is cut by a time-dependent level-set function. By construction, cut mesh elements are general curved space-time elements, is allowed by the analyzed DG method.
This approach, similar to cut-FEM methods, avoids the need to remesh the computational domain at every time step. In this setting, agglomeration plays a closely related but different role from its use in the construction of coarse polygonal grid hierarchies: background cells whose intersection with the physical domain is relatively small are merged with suitable neighbouring elements to form robust polytopic space--time elements.
Under the geometric assumptions used in the analysis, the resulting polytopic space--time elements satisfy bounds that are uniform with respect to the cut configuration. This work therefore extends the agglomeration concepts developed in Work Package 1.3 to unfitted, time-dependent geometries. A preliminary numerical implementation has also been developed to illustrate the resulting mesh construction and discrete solution (Fig. \ref{fig:sissa_mesh}).

\begin{figure}
    \begin{subfigure}[h]{0.48\textwidth}
        \centering
        \includegraphics[width=1\textwidth]{figures/bg_mesh.png}
        \caption{}
        \label{bg_mesh}
    \end{subfigure}
    \hfill
    \begin{subfigure}[h]{0.48\textwidth}
        \centering
        \includegraphics[width=1\textwidth]{figures/cut_mesh.png}
        \caption{}
        \label{cut_mesh}
    \end{subfigure}
    \centering
    \caption{Illustration of the unfitted space–time discretization. Left: the background space–time mesh before restriction to the physical domain. Right: the active cut mesh and a representative numerical solution on the time-dependent physical domain.}
    \label{fig:sissa_mesh}
\end{figure}

A space–time DG formulation has been developed on the resulting agglomerated mesh. Stability has been established through coercivity and inf-sup estimates that are uniform with respect to the position of the moving boundary relative to the background mesh. An a priori error estimate has also been derived in a DG energy norm whose spatial component controls the broken $H^1$-seminorm. For fixed polynomial degrees, the estimate is optimal with respect to both the spatial mesh size and the time-step size. The corresponding $L^2$-error analysis has not yet been completed and constitutes the next theoretical step.



\newpage

\section{WP1.4: Integration of PSCToolkit within deal.II (UNITOV)}
\label{sec:unitov}

\subsection{Integration status of PSCToolkit}
The integration of the capabilities provided by PSCToolkit into the \texttt{deal.II} repository has moved
towards a more mature stage. The numerical linear algebra backbone of PSCToolkit, provided by the PSBLAS library, has
at this point been fully integrated, mainly through PRs~\href{https://github.com/dealii/dealii/pull/18938}{\#18938} and~\href{https://github.com/dealii/dealii/pull/19356}{\#19356}.\footnote{links retrieved on 30.07.2026} The design of the user interface
closely follows the existing structure of the main \texttt{deal.II} linear algebra frameworks, such as the
interfaces to Trilinos and PETSc. In particular, the namespace \texttt{PSCToolkitWrappers} defines the
interface to the PSCToolkit classes. In addition, several unit tests have been added to ensure the correctness of the implementation
and to avoid regressions in future developments.

In the same fashion, AMG preconditioners are provided through AMG4PSBLAS, a dedicated package
of algebraic multilevel preconditioners tailored for PSBLAS data structures. Its integration is
carried out in PR~\href{https://github.com/dealii/dealii/pull/19398}{\#19398},\footnote{retrieved on 30.07.2026} and is currently under review.

This interface is realized by the class \texttt{PreconditionAMG}, which provides
a minimal set of methods for constructing and applying AMG preconditioners. The action of
the (inverse of the) preconditioner on a vector is realized through the standardized \texttt{vmult()} method. The preconditioner is easily customizable through an \texttt{AdditionalData} object, which allows specifying the
type of smoothers, the coarsening strategy, and the coarse grid solver. An example of the usage of the preconditioner from AMG4PSBLAS is shown in Listing~\ref{lst:linear_solve}.
We refer to the AMG4PSBLAS documentation for a complete and detailed description of the available options.

\begin{lstlisting}[language=C++,caption={Preconditioned conjugate gradient solver class with PSCToolkit wrappers.}, label=lst:linear_solve]
void MyProblem::solve()
{
    PSCToolkitWrappers::PreconditionAMG::AdditionalData prec_data;
    prec_data.cycle_type       = "VCYCLE";
    prec_data.smoother_type    = "POLY";
    prec_data.aggregation_type = "MATCHBOXP";

    PSCToolkitWrappers::PreconditionAMG preconditioner;
    preconditioner.initialize(system_matrix, prec_data);
    solver.solve(system_matrix,
                 solution,
                 system_rhs,
                 preconditioner);
}
\end{lstlisting}

\subsection{Results on MareNostrum 5}

The interface has first been tested by developing a simple tutorial application (an early version is
documented in deliverable D1.6), which shows the seamless integration of the PSCToolkit interface into \texttt{deal.II}. Having validated the implementation, we have
started assessing the practical performance of the AMG preconditioners on the MareNostrum 5 supercomputer
at the Barcelona Supercomputing Center\footnote{MareNostrum 5 information: \url{https://www.bsc.es/marenostrum/marenostrum-5}, retrieved on 30.07.2026}.

As an initial benchmark, we consider a three-dimensional Poisson problem on the unit cube, discretized
with $\mathcal{Q}_1$ finite elements on a uniform hexahedral mesh. The resulting linear system is solved with a conjugate
gradient method preconditioned by a single V-cycle of AMG provided by AMG4PSBLAS. The preconditioner is configured
as follows:
\begin{itemize}
    \item \textbf{Cycle}: smoothed V-cycle with parallel matching aggregation (\texttt{MatchBoxP}, $3$ sweeps).
    \item \textbf{Smoother}: polynomial smoother (\texttt{CHEB\_1\_OPT}) of degree $4$, with a local $\ell_1$ diagonal solver.
    \item \textbf{Coarsest level}: $\ell_1$-Jacobi solver ($30$ sweeps).
\end{itemize}


The aggregation strategy is described in~\cite{BootCMatch}. We report weak-scaling results obtained on 1, 8, and 64 nodes (24 MPI ranks per node), corresponding to
approximately $2.15\times 10^6$, $1.70\times 10^7$, and $1.35\times 10^8$ unknowns, respectively.
Figure~\ref{fig:psctoolkit_weak_scaling} summarizes two key indicators extracted from the runs:
the average time per CG iteration and the operator complexity of the AMG hierarchy.
The iteration cost remains in a narrow range (about $0.08$-$0.10$ s/iteration), while operator complexity
stays constant at almost $2$, indicating that the hierarchy quality is preserved as the problem size grows.

\begin{figure}[htbp]
    \centering
    \begin{subfigure}[t]{0.49\textwidth}
        \centering
        \includegraphics[width=\textwidth]{figures/weak_scaling_time_per_iter.pdf}
    \end{subfigure}
    \hfill
    \begin{subfigure}[t]{0.49\textwidth}
        \centering
        \includegraphics[width=\textwidth]{figures/weak_scaling_coarsening.pdf}
    \end{subfigure}
    \caption{Weak-scaling behavior of AMG4PSBLAS on MareNostrum 5. Left: average time per CG iteration. Right: AMG operator complexity.}
    \label{fig:psctoolkit_weak_scaling}
\end{figure}

\begin{figure}[htbp]
	\centering
	\begin{subfigure}[t]{0.32\textwidth}
		\centering
		\includegraphics[width=\textwidth]{figures/left_ventricle_view.pdf}
		\caption{Left ventricle mesh.}
		\label{fig:left_ventricle_mesh}
	\end{subfigure}
	\hfill
	\begin{subfigure}[t]{0.55\textwidth}
		\centering
		\includegraphics[width=\textwidth]{figures/strong_scaling_ventricle.png}
		\caption{Strong-scaling analysis on MN5.}
		\label{fig:left_ventricle_strong_scaling_plot}
	\end{subfigure}
	\caption{Strong-scaling study of AMG4PSBLAS on an unstructured left ventricle mesh ($\approx 2.4\times 10^7$ DoFs) for the monodomain model in cardiac
		electrophysiology. Left: the computational mesh. Right: solve time as a function of the
		number of MPI ranks.}
	\label{fig:left_ventricle_strong_scaling}
\end{figure}

To assess the performance of the AMG preconditioners with a more realistic test case, we
consider the linear system arising from the monodomain model in cardiac electrophysiology.
The problem is discretized on a fixed, unstructured, patient-specific left ventricle mesh
(shown in Figure~\ref{fig:left_ventricle_strong_scaling}) with approximately $2.4\times 10^7$
degrees of freedom. The AMG4PSBLAS preconditioner is configured as in the previous Poisson test case, and the linear
system is solved with a preconditioned conjugate gradient method, since the system matrix is symmetric and positive definite.
With this configuration, the conjugate gradient solver reduces the relative residual by eight orders of magnitude
in $12$ iterations, independently of the number of MPI ranks employed. Figure~\ref{fig:left_ventricle_strong_scaling}
reports the strong-scaling behavior.



\subsection{Ongoing and future work}

Ongoing work aims at finalizing the integration of GPU support and at assessing the performance of the AMG4PSBLAS preconditioners in practical
scenarios. A first line of work concerns GPU support: PSBLAS provides GPU acceleration mainly through CUDA backends, and the exposition
of this functionality through the \texttt{deal.II} interface has already been implemented, with a
corresponding pull request to be submitted in the near future. In parallel, we plan to carry out larger-scale performance
evaluations on HPC systems and to test the effectiveness of the AMG preconditioners on coupled multiphysics problems. Finally, in collaboration with
INPT, we are exploring the use of Block Low-Rank approximations as smoothers within the AMG cycles.



\newpage



\section{WP1.5: Integration of MUMPS within deal.II (INPT)}
\label{sec:inpt}

%\lipsum[8-9]
In this section, we highlight some of the recent work carried out with
the MUMPS direct solver in the context of the dealii-X project. The
primary objective has been to evaluate the use of low rank
approximation techniques and mixed precision algorithms for organ
simulation problems arising from \texttt{deal.II}. The following
sections provide a brief background on the theoretical aspect and the
recent features available in MUMPS with performance results on a
sample matrix. The sample matrix is a medium resolution matrix for the
brain indentation problem obtained from our collaborators at FAU
Erlangen. The problem consists of 3 million degrees of freedom and
$1.2\times 10^3$ elements. We solve the problem with the direct solver
for a desired accuracy of $10^{-10}$. All the results in this section
have been obtained on 1 compute node with 24MPIx8th configuration on
the CPU partition of the Kairos supercomputer in the CALMIP supercomputing
center in Toulouse (France).

\subsection{Low Rank Approximation}
% \begin{itemize}[left=1em, itemsep=0pt, topsep=0pt] 
	%     \item \lipsum[10][1-2]
	%     \item \lipsum[10][3-4]
	%     \item \lipsum[10][5-6]
	% \end{itemize}
Data sparsity within a matrix results in a low numerical rank. We can
take advantage of this fact by using a low-rank approximation for the
respective matrix. MUMPS uses the Block Low-Rank approach to
represent low rank blocks within the dense frontal matrices during the
numerical factorization phase~\citep{a.b.h.l.:22,a.b.l.m:17}. This can
result in a significant reduction in memory consumption and
consequently lower operation count. The accuracy of the low rank
representation is controlled by a user-defined parameter
$\varepsilon$, known as the BLR threshold,
\begin{equation}\label{BLR_error}
	\left\lVert A-X\Sigma Y^T\right\rVert_2 \leq \varepsilon,
\end{equation}
where $X\Sigma Y^T$ represents the truncated SVD of $A$.
Additionally, MUMPS is capable of storing different columns of the
factorized block, i.e. $X$ and $Y$ in the above example, at different
precisions, based on the magnitude of the corresponding singular
value. This is referred to as adaptive precision BLR~\citep{a.b.b.g:22}
and leads to a further reduction in memory footprint.

\begin{figure}
	\begin{subfigure}[h]{0.48\textwidth}
		\centering
		\includegraphics[width=1\textwidth]{brain_LU_size}
		\caption{}
		\label{LU_size}
	\end{subfigure}
	\hfill
	\begin{subfigure}[h]{0.48\textwidth}
		\centering
		\includegraphics[width=1\textwidth]{brain_peak_mem}
		\caption{}
		\label{peak_mem}
	\end{subfigure}
	\centering
	\caption{Memory study: (a) size of $LU$ factors and (b) peak memory usage during factorization. The bar graphs show the comparison between double precision and adaptive precision BLR for different values of the BLR threshold $\varepsilon$. FR refers to the full rank case.}
	\label{fig:brain_mem}
\end{figure}

\begin{figure}
	\begin{subfigure}[h]{0.48\textwidth}
		\centering
		\includegraphics[width=1\textwidth]{figures/brain_flops}
		\caption{}
		\label{flops}
	\end{subfigure}
	\hfill
	\begin{subfigure}[h]{0.48\textwidth}
		\centering
		\includegraphics[width=1\textwidth]{figures/brain_t_fac}
		\caption{}
		\label{t_fac}
	\end{subfigure}
	\centering
	\caption{(a) Operation count and (b) elapsed time during the numerical factorization phase. The bar graphs show the comparison between double precision and adaptive precision BLR for different values of the BLR threshold $\varepsilon$. FR refers to the full rank case.}
	\label{fig:brain_fac}
\end{figure}

Figure \ref{fig:brain_mem} illustrates the memory footprint for the sparse direct solver for different values of the BLR threshold $\varepsilon$. We can observe that with the introduction of BLR, the size of the $LU$ factors reduces by $28\%$. The gain goes up to $50\%$ for the largest possible $\varepsilon$ of $10^{-7}$. Using adaptive precision BLR further brings down the memory consumption by around $13\%$. Similarly, the peak memory used during the factorization process reduces by $26\%$ for the largest BLR threshold compared to the full rank case.

Figure \ref{fig:brain_fac} highlights the operation counts and factorization time for solving the brain matrix problem. The FLOPS counts reduces drastically by $60\%$ owing to BLR compression. It goes down consistently and achieves a gain of $85\%$ for the largest $\varepsilon$ value. The time elapsed during factorization reduces by $32\%$ for $\varepsilon=10^{-7}$ compared to the full rank case. These results clearly illustrate the benefits of using low rank approximation techniques both in terms of reducing complexity and improving performance. This makes the use of MUMPS possible even for memory intensive applications.

\subsection{Mixed Precision Iterative Refinement}
%\lipsum[11]
Mixed precision algorithms refer to the use of multiple floating point formats with different precisions within a given algorithm. In this work, we will focus on mixed LU-IR where we perform certain steps of the iterative refinement algorithm at a lower precision. The goal of iterative refinement is to recover the loss in accuracy due to approximate computing. Algorithm \ref{algorithm_1} lays out the procedure for performing mixed LU-IR with two precisions: single (fp32) and double (fp64). The idea is to perform the time and memory intensive steps at a lower precision in order to reduce the overall memory footprint and the execution time of the algorithm.
\begin{algorithm}
	\caption{Mixed LU-IR}
	\label{algorithm_1}
	\begin{algorithmic}[1]
		\STATE Factorize $A=LU$ \hfill (in precision fp32)
		\STATE Solve $Ax_1=b$ via $x_1=U^{-1}(L^{-1}b)$ \hfill (in precision fp32)
		\REPEAT  
		\STATE \quad $r_i = b - Ax_i$ \hfill (in precision fp64)
		\STATE \quad Solve $Ad_i=r_i$ via $d_i=U^{-1}(L^{-1}r_i)$ \hfill (in precision fp32)
		\STATE \quad $x_{i+1} = x_i + d_i$ \hfill (in precision fp64)\;
		\UNTIL target tolerance reached
	\end{algorithmic}
\end{algorithm}

The above algorithm results in $\mathcal{O}(n^3)$ operations in single
precision and $\mathcal{O}(n^2)$ operations in double precision per
iteration. It is guaranteed to converge in double precision given that
$\kappa(A) < \frac{1}{u_{\text{fp32}}}=1.6 \times 10^7$, where
$\kappa(A)$ is the condition number of the matrix $A$ and
$u_{\text{fp32}}$ is the unit round-off for single precision
float~\citep{a.b.h.l.:22,c.h:18}.

\begin{table}[htbp]
	\centering
	\caption{Performance comparison between double precision and single precision factorization with iterative refinement.}
	\begin{center}
		\begin{tabular}{ p{1.5cm}| p{1.5cm}| p{1cm}| p{1.7cm}| p{1cm}| p{1cm}| p{1cm}| p{1.2cm}| p{1.2cm} }
			%\begin{tabular}{ c|c|c|c|c|c|c|c|c }
			% \hline
			\hline
			\hline
			\multirow{2}{2em}{}    & \multicolumn{2}{c|}{Memory study (GB)} & \multirow{2}{2em}{FLOPS} & \multicolumn{3}{c|}{Elapsed time (sec)} & \multicolumn{2}{c}{Bwd error}                                            \\
			& Size $\vert LU\vert$                   & Memory used              &                                         & Facto                         & Solve & LU-IR (nb) & Initial  & LU-IR    \\
			\hline
			\hline
			Full rank (FR)         & 137.04                                 & 221.91                   & 2.23E+14                                & 35.55                         & 0.21  & $-$        & 2.81E-16 & $-$      \\
			\hline
			BLR $\varepsilon$=1E-7 & 64.88                                  & 163.68                   & 3.33E+13                                & 24.17                         & 0.12  & 8.19 (3)   & 3.35E-6  & 1.94E-13 \\
			\hline
			\multicolumn{9}{c}{Single precision factorization}                                                                                                                                                              \\
			\hline
			Full rank (FR)         & 68.52                                  & 112.36                   & 2.23E+14                                & 18.43                         & 0.12  & 8.08 (3)   & 5.11E-11 & 1.39E-12 \\
			\hline
			BLR $\varepsilon$=1E-7 & 32.44                                  & 84.06                    & 3.33E+13                                & 13.08                         & 0.08  & 9.37 (4)   & 2.84E-6  & 1.91E-14 \\
			\hline
			\hline
		\end{tabular}
	\end{center}
	\label{table:SinD-IR}
\end{table}

Table \ref{table:SinD-IR} captures the different quantities of interest and compares them for double precision and single precision factorization with iterative refinement.
%Single in double refers to performing single precision factorization within a double precision instance of MUMPS (algorithm \ref{algorithm_1}). 
We are able to recover the loss in accuracy by employing the iterative refinement technique. We can observe that by utilizing the mixed LU-IR approach of algorithm \ref{algorithm_1}, we are able to reduce the memory consumption and factorization time by nearly half. Thus, by combining mixed LU-IR with BLR compression, we can bring down memory consumption by $77\%$ and factorization time by $63\%$. This will allow us to use fewer computational resources even for larger problems.

We plan to further investigate the use of GPU acceleration features of MUMPS to factorize the dense frontal matrices. We are also working on the post-processing BLR feature, which allows us to compress the low rank blocks on the CPU after the full rank factorization has been carried out on the GPU. Finally, we plan to extend the study to more organ simulations such as problems involving the liver or the lungs.

\newpage

\section{{Conclusion}} \label{sec:conclusion}

This report has documented the improvements made to \texttt{deal.II} during the fourth
semester of the dealii-X project. The RUB group extended the matrix-free GPU
infrastructure with a batched processing strategy that removes the thread-block
occupancy limit for low-order elements, doubling throughput at polynomial order one on
the NVIDIA GH200, and applied the same techniques to Raviart--Thomas elements, to
discontinuous Galerkin operators with face integration, and to a portable multigrid
solver; the matrix-free transfer operators have been merged into
\texttt{deal.II} and will appear in release 9.8. UNIPI completed an implicit HDG solver
for one-dimensional blood flow, coupled monolithically to the \textsc{SUNDIALS} IDA
time integrator, and validated it against the 37-artery systemic network benchmark.
SISSA saw the agglomeration-based polygonal Poisson program merged into the
\texttt{deal.II} Code Gallery and submitted the reusable components of \texttt{Polydeal}
to the core library, while extending the same agglomeration ideas to unfitted
space-time discretizations on moving domains. UNITOV completed the integration of the
PSBLAS linear algebra backbone, opened the interface to the AMG4PSBLAS preconditioners
for review, and demonstrated stable weak and strong scaling on MareNostrum 5, including
a patient-specific left ventricle monodomain problem. INPT showed that block low-rank
compression combined with mixed-precision iterative refinement reduces the memory
footprint of the MUMPS factorization of a brain indentation problem by up to $77\%$
and the factorization time by $63\%$, without loss of accuracy.

Taken together, these results indicate that the components developed in the earlier
semesters are maturing into production functionality: the majority of the work reported
here is either already part of \texttt{deal.II} or under review for inclusion. The
activities planned for the remainder of the project follow directly from the open
points identified in the individual sections, namely the reduction of communication
volume and a full scalability study for the DG multigrid solvers, the completion of the
$L^2$-error analysis for the unfitted space-time method, the finalisation of GPU
support in the PSCToolkit interface and its assessment on coupled multiphysics problems,
and the extension of the low-rank and mixed-precision studies to GPU factorization and
to further organ simulations.

\bibliographystyle{plainnat}
\bibliography{../common/bibliography}

\label{MyLastPage}

\end{document}

%%% Local Variables:
%%% mode: LaTeX
%%% TeX-master: t
%%% End:
